% =============================================================================
% Section 1: Introduction
% =============================================================================

\section{Introduction}
\label{sec:introduction}

% AGENT: Fill physics motivation from theory scout Wave 1 report
% Describe the theoretical context, why this search is important,
% and what gap in knowledge it addresses.

The Standard Model of particle physics has been remarkably successful in
describing fundamental interactions. However, several open questions remain,
including [theoretical motivation]. This analysis searches for [signal process]
which would provide evidence for [physics beyond the Standard Model / rare
process measurement].

% AGENT: Fill signal process description from analysis strategy document
% Include Feynman diagram reference, production mechanism, decay chain,
% and expected final state signature.

\subsection{Signal Process}
\label{sec:intro:signal}

The signal process under study is [process name], produced via [production
mechanism] and decaying to [final state]. The expected cross-section at
\sqrts = [energy] is [cross-section value] as computed at [order] in
perturbative QCD.

% AGENT: Fill previous measurements from literature review
% Include references to prior searches, existing limits, and how this
% analysis improves upon them.

\subsection{Previous Measurements}
\label{sec:intro:previous}

Previous searches for this process have been performed by [experiments].
The most stringent existing limit is [limit value] at 95\% confidence level,
obtained by [experiment] using [dataset]~\cite{previous_result}.
This analysis extends the sensitivity by [improvement description].

% AGENT: Fill analysis overview summarizing the strategy
% Brief roadmap of the note structure.

\subsection{Analysis Overview}
\label{sec:intro:overview}

This note describes a search for [signal] using data collected by [experiment]
at \sqrts = [energy], corresponding to an integrated luminosity of \lumival.
The analysis employs [key technique, e.g., BDT-based signal extraction] with
[background estimation method]. The statistical interpretation uses a
profile likelihood ratio test implemented in \texttt{pyhf}.

The note is organized as follows: Section~\ref{sec:data_mc} describes the
data and Monte Carlo samples; Section~\ref{sec:objects} defines the
reconstructed objects; Section~\ref{sec:selection} presents the event
selection; Section~\ref{sec:background} details the background estimation;
Section~\ref{sec:systematics} discusses systematic uncertainties;
Section~\ref{sec:statistical} describes the statistical analysis;
and Section~\ref{sec:expected} presents the expected results.
